\section{简单即美：拆掉中间层，也拆掉组织层级}

前面讲的是技术层面的拆，本章转到组织层面的拆。刘先明的“简单即美”不是审美口号，而是工程管理原则：复杂系统要想迭代快，就必须减少不必要的中间层。

本章从技术上的“拆”转向组织上的“简化”。刘先明说自己喜欢简单，上任后最重要的决策就是简化流程、简化研发工序、合并重复事项、降低不必要事项优先级。他不喜欢只听汇报，而是直接看前线问题、代码和实验结果。

\begin{figure}[H]
\centering
\includegraphics[width=0.96\textwidth]{simple-is-beautiful.png}
\caption{简单即美：不断拆掉浓郁中间层，让数据直接训练能力。自制概念图，依据 00:00:00--00:02:16 与 00:54:30--01:48:46 对谈内容整理。}
\end{figure}

\begin{knowledgebox}{读图：技术简化和组织简化是同一件事}
拆激光雷达依赖、拆规控规则、拆中间层、拆 Language，都是为了减少瓶颈；组织上减少层级、直接看问题、合并重复工作，也是为了减少信息瓶颈。
\end{knowledgebox}

\subsection{扁平工程组织}

在物理 AI 和量产结合的场景里，决策速度很重要。刘先明希望层级被拍平，一线 engineer 能直接暴露问题，团队 lead 也能直接写代码、看实验、做 deep dive。组织内部不设置太多部门墙，避免 duplicate work，通过共享实验和失败来提高整体速度。

\begin{figure}[H]
\centering
\includegraphics[width=0.96\textwidth]{org-flat-engineering.png}
\caption{扁平工程组织：前沿 AI 与量产结合，需要一线问题快速决策。自制概念图，依据 00:54:30--01:48:46 对谈内容整理。}
\end{figure}

\begin{knowledgebox}{读图：扁平不是无管理，而是短链路}
一线工程师暴露代码和实验，负责人直接看问题，deep dive 复盘好坏实验，资源整合减少重复，最终支持模型超越规则后的快速切换。
\end{knowledgebox}

\subsection{工程文化：Deep Dive、共享失败和减少重复}

本节把“扁平”进一步落成日常机制。没有 deep dive 和失败共享，扁平组织很容易只是少几层汇报线，而不是真正提高学习速度。

访谈中提到团队会做 deep dive，把好的实验、坏的实验、最近看到的东西、自己做挂的东西拿出来分享。这个做法服务于两个目标：一是让团队快速形成共同上下文，二是避免不同小组重复造轮子。Physical AI 的研发需要大量实验，如果失败不能共享，组织会反复踩同一个坑。

\begin{importantbox}{工程组织的知识复用}
前沿 AI 团队的核心资产不只是模型权重，还包括失败实验、调参经验、数据处理方法和对问题的共同理解。Deep Dive 是把这些隐性知识显性化。
\end{importantbox}

\subsection{切换时机：机会与风险}

接下来讨论简化最危险的一步：什么时候可以把旧路线切掉。对车企来说，这不是研究组内部开关，而是直接影响用户安全和量产质量的决策。

车企不能因为技术酷就上线。模型要真正替代规则，需要等到某个时间点：模型性能明显超过规则，坏处可控，好处明显，安全责任和量产质量能承受。刘先明提到测试员鼓掌的时刻，就是模型超越旧规则的一种组织信号。

\begin{figure}[H]
\centering
\includegraphics[width=0.92\textwidth]{ai-transition-risk-balance.png}
\caption{切换时机：机会 vs 风险。车企必须等模型超过规则，而不是只因技术酷就上线。}
\end{figure}

\begin{warningbox}{技术路线切换不能靠信仰硬推}
在车企里，技术路线切换必须对安全、质量和用户负责。正确做法不是“模型很酷所以全量上线”，而是让模型在关键指标上超过旧系统，并逐步扩大部署。
\end{warningbox}

\subsection{本章小结}

小鹏这次转型同时发生在技术和组织两层：技术上拆掉中间层，组织上缩短信息链路。两者都服务于同一个目标：让数据和真实问题更快进入模型迭代。

